\clearpage
\section{Intrinsic multiblock
geometry}\label{a-intrinsic-multiblock-geometry}

Technical derivation for review. Equation and subsection numbers in this
appendix are local to the source derivation. Historical local labels are
retained, including numbering gaps or nonsequential labels. The
accompanying source notes retain the original expressions for
comparison.

\subsection{1. Setting and precise
scope}\label{a-setting-and-precise-scope}

Fix \(N\) and a partition \(N=n_{1}+\ldots +n_{k}, k\ge 2\). Hermitian
matrix tuples have the real trace inner product. Retain the residual
group \(K=S(U(n_{1}) \times  \ldots  \times  U(n_{k}))\), including its
action on the fermions; do not divide by any internal nonabelian orbit.
Write

% DISPLAY a.1
% SOURCE X_i=Z_i+A_i+Y_i,
% SOURCE Z_i=diag(z_1i I_(n_1),...,z_ki I_(n_k)),
% SOURCE sum_a n_a z_a=0,   Tr A_i^(a)=0,
% SOURCE b_ab=z_a-z_b,   r_ab=|b_ab|,   r_*=min_(a<b) r_ab>0.
\begin{gather*}
\begin{aligned}
& X_{i}=Z_{i}+A_{i}+Y_{i}, \\
& Z_{i}=\operatorname{diag} (z_{1i} I_{n_{1}},\ldots ,z_{ki} I_{n_{k}}), \\
& \sum _{a} n_{a} z_{a}=0, \quad  \operatorname{Tr}  A_{i}^{(a)}=0, \\
& b_{ab}=z_{a}-z_{b}, \quad  r_{ab}=\vert b_{ab}\vert , \quad  r_{*}=\min _{a<b} r_{ab}>0.
\end{aligned}\notag
\end{gather*}

A is block diagonal and \(Y\) is off block. Put
\(\alpha _{a}^{2}=\sum _{i} \Vert A_{i}^{(a)}\Vert _{F}^{2}\). Assume
only

% DISPLAY a.2
% SOURCE alpha_a+alpha_b <= kappa r_ab  for every a<b.             (1)
\begin{gather*}
\begin{aligned}
& \alpha _{a}+\alpha _{b} \le  \kappa  r_{ab} \quad  \text{ for every } a<b.
\end{aligned}\tag{1}
\end{gather*}

There is no hypothesis \(\alpha _{a} \ll  r_{*}\) for a block unrelated
to a shortest pair. Set \(\kappa \le 1/100\). For the fast tube use

% DISPLAY a.3
% SOURCE |Y| <= sigma r_*,   sigma<=1/100.                         (2)
\begin{gather*}
\begin{aligned}
& \vert Y\vert  \le  \sigma  r_{*}, \quad  \sigma \le 1/100.
\end{aligned}\tag{2}
\end{gather*}

Condition (2) is deliberately a fast-coordinate condition. It is
compatible with the Gaussian widths at sufficiently large \(r_{*}\) and
does not impose the forbidden common scale on A. A smooth smaller cutoff
may use the harmonic minimum of the \(r_{ab}\) instead of the nonsmooth
literal minimum.

Let \(C\) be the Euclidean block-center tuple space, I the space of
blockwise traceless internal tuples, and \(m\) the off-block Hermitian
gauge-parameter space. Its real dimension is

% DISPLAY a.4
% SOURCE q=2 sum_(a<b) n_a n_b.
\begin{gather*}
\begin{aligned}
& q=2 \sum _{a<b} n_{a} n_{b}.
\end{aligned}\notag
\end{gather*}

For a Hermitian tuple \(Q\) let \(B_{Q} \xi =(i[\xi ,Q_{i}])_{i}\).
Define

% DISPLAY a.5
% SOURCE R=(B_Z^* B_Z)^(1/2)=direct_sum_(a<b) r_ab I_(2n_an_b),
% SOURCE E=B_Z R^(-1),   E^*E=I_m,
% SOURCE N_Z=ker B_Z^* in the off-block tuple space.
\begin{gather*}
\begin{aligned}
& R=(B_{Z}^{*} B_{Z})^{1/2}=\bigoplus _{a<b} r_{ab} I_{2n_{a} n_{b}}, \\
& E=B_{Z} R^{-1}, \quad  E^{*}E=I_{m}, \\
& N_{Z}=\ker  B_{Z}^{*} \text{ in the off-block tuple space }.
\end{aligned}\notag
\end{gather*}

The intrinsic normal constraints are \(B_{Z}^{*}Y=0\), equivalently
\(b_{ab} \cdot  Y_{ab}=0\) for every pair. The slice \(S\) consists of
the triples \((Z,A,Y)\) satisfying these constraints. Statements below
are tensor identities on this slice and its local gauge-saturated image.
A simultaneous smooth projector branch exists locally by the nonsingular
Hessian proved below. This does not assert a globally unique branch
across incompatible cluster partitions.

\subsection{2. Horizontal differential and the exact internal
coefficient}\label{a-horizontal-differential-and-the-exact-internal-coefficient}

Use the normal connection obtained by orthogonal projection onto
\(N_{Z}\). For \(h=(\delta  Z,\delta  A,\delta  Y_{N})\) in
\(H=C \oplus  I \oplus  N_{Z}\), horizontal differentiation of
\(B_{Z}^{*}Y=0\) gives

% DISPLAY a.6
% SOURCE delta Y=delta Y_N+E K h,
% SOURCE K h= -R^(-1) B_(delta Z)^*Y.                             (3)
\begin{gather*}
\begin{aligned}
& \delta  Y=\delta  Y_{N}+E K h, \\
& K h= -R^{-1} B_{\delta  Z}^{*}Y.
\end{aligned}\tag{3}
\end{gather*}

\(K\) vanishes identically on I and \(N_{Z}\). The slice embedding
differential, in the orthogonal ambient splitting
\(H \oplus  \operatorname{range}  E\), is

% DISPLAY a.7
% SOURCE S h=(h,K h),
% SOURCE g_S=S^*S=I_H+K^*K
% SOURCE    =g_C direct_sum I_I direct_sum I_(N_Z).                (4)
\begin{gather*}
\begin{aligned}
& S h=(h,K h), \\
& g_{S}=S^{*}S=I_{H}+K^{*}K \\
& =g_{C} \oplus  I_{I} \oplus  I_{N_{Z}}.
\end{aligned}\tag{4}
\end{gather*}

Thus the induced slice has internal kinetic coefficients \textbf{exactly
one}, and there are no internal/center or internal/fast induced-metric
cross terms. The full reduced co-metric additionally contains the
positive orbital contribution from the Gauss square; it is not claimed
that its total AA block equals I. The coefficient-one \(t_{A}\) summand
is retained before that square is estimated. This is an exact identity
for every A and \(Y\) in the coordinate domain, not an oscillator
expectation.

Skew-adjointness of the gauge action gives a useful equivalent identity.
For \(O=B_{X}, O_{C}=P_{C} O\) and \(O_{H}=P_{H} O\),

% DISPLAY a.8
% SOURCE O_C=P_C B_Y,   K=R^(-1) O_C^*.                            (5)
\begin{gather*}
\begin{aligned}
& O_{C}=P_{C} B_{Y}, \quad  K=R^{-1} O_{C}^{*}.
\end{aligned}\tag{5}
\end{gather*}

The elementary commutator bound \(\Vert B_{Y}\Vert \le 2\vert Y\vert\)
proves

% DISPLAY a.9
% SOURCE ||K||<=2s,   s=|Y|/r_*,
% SOURCE (1+4s^2)^(-1) I_C <= g_C^(-1) <= I_C.                    (6)
\begin{gather*}
\begin{aligned}
& \Vert K\Vert \le 2s, \quad  s=\vert Y\vert /r_{*}, \\
& (1+4s^{2})^{-1} I_{C} \le  g_{C}^{-1} \le  I_{C}.
\end{aligned}\tag{6}
\end{gather*}

All normal-frame connections remain inside the covariant center
momentum. Their unbounded \(Y \partial _{Y}\) generators are never
estimated by oscillator energy on arbitrary vectors.

\subsection{3. Faddeev-Popov matrix, explicit noncircular invertibility,
and
density}\label{a-faddeev-popov-matrix-explicit-noncircular-invertibility-and-density}

The differential of the normalized slice constraint is
\(N^{*}=(-K,I_{m})\); thus

% DISPLAY a.10
% SOURCE N=(-K^*,I_m),   W=N^*N=I_m+K K^*.
\begin{gather*}
\begin{aligned}
& N=(-K^{*},I_{m}), \quad  W=N^{*}N=I_{m}+K K^{*}.
\end{aligned}\notag
\end{gather*}

The full transverse-orbit matrix is

% DISPLAY a.11
% SOURCE M=N^*O=E^*B_X-K O_C
% SOURCE  =E^*B_(Z+A)+E^*B_Y-R^(-1)O_C^*O_C.                     (7)
\begin{gather*}
\begin{aligned}
& M=N^{*}O=E^{*}B_{X}-K O_{C} \\
& =E^{*}B_{Z+A}+E^{*}B_{Y}-R^{-1}O_{C}^{*}O_{C}.
\end{aligned}\tag{7}
\end{gather*}

Put \(F=E^{*}B_{Z+A}\). \(F\) is pair-block diagonal. Under the
canonical complex identification of a pair, its block is unitarily
equivalent to

% DISPLAY a.12
% SOURCE r_ab I+T_(ab,n),
% SOURCE T_(ab,n) V=A_n^(a)V-V A_n^(b),  n=b_ab/r_ab.
\begin{gather*}
\begin{aligned}
& r_{ab} I+T_{ab,n}, \\
& T_{ab,n} V=A_{n}^{(a)}V-V A_{n}^{(b)}, \quad  n=b_{ab}/r_{ab}.
\end{aligned}\notag
\end{gather*}

Consequently

% DISPLAY a.13
% SOURCE (1-kappa)R <= F <= (1+kappa)R,   FR=RF.                  (8)
\begin{gather*}
\begin{aligned}
& (1-\kappa )R \le  F \le  (1+\kappa )R, \quad  FR=RF.
\end{aligned}\tag{8}
\end{gather*}

The transverse-frame convention \(E=B_{Z} R^{-1}\) has absorbed the
fixed real root-plane complex structure; that is why \(F\) is positive
here. In a fixed real root frame the corresponding matrix carries that
orthogonal complex structure.

Equations (5)-(8) give the explicit bound

% DISPLAY a.14
% SOURCE ||M R^(-1)-I|| <= kappa+2s+4s^2 <=0.0304.                (9)
\begin{gather*}
\begin{aligned}
& \Vert M R^{-1}-I\Vert  \le  \kappa +2s+4s^{2} \le 0.0304.
\end{aligned}\tag{9}
\end{gather*}

No \(r_{\max}/r_{*}\) or unrelated internal radius enters (9). The full
differential from \((h,\xi )\) to ambient coordinates has block matrix

% DISPLAY a.15
% SOURCE T=[ I_H   O_H ]
% SOURCE   [ K     E^*O].
\[T=\begin{bmatrix}I_H&O_H\\K&E^{*}O\end{bmatrix}.\]

Its determinant is det \(M\). Equivalently the induced slice volume is
\(\sqrt{\det  W}\), while the normalized-normal orbit contraction is
\(\vert \det (W^{-1/2}M)\vert\). They cancel. The physical measure after
the compact off-block gauge integration is therefore

% DISPLAY a.16
% SOURCE J dZ dA dY_N,   J=det M>0,                              (10)
\begin{gather*}
\begin{aligned}
& J dZ dA dY_{N}, \quad  J=\det  M>0,
\end{aligned}\tag{10}
\end{gather*}

up to one constant orbit normalization. This is physical measure
including orbit volume, not bare quotient Riemannian volume.

At \(Y=0\),

% DISPLAY a.17
% SOURCE J_0=det F=product_(a<b) det_C(r_ab I+T_(ab,n))^2,
% SOURCE d_0=sqrt(J_0)=product_(a<b) det_C(r_ab I+T_(ab,n)).        (11)
\begin{gather*}
\begin{aligned}
& J_{0}=\det  F=\prod _{a<b} \det _{C}(r_{ab} I+T_{ab,n})^{2}, \\
& d_{0}=\sqrt{J_{0}}=\prod _{a<b} \det _{C}(r_{ab} I+T_{ab,n}).
\end{aligned}\tag{11}
\end{gather*}

The determinant is positive by (1). Internal coincidences introduce no
singularity.

There is no linear-Y density term, even for nonzero internal A. Indeed
\(L=E^{*}B_{Y}\) has zero diagonal pair blocks: acting on an ab gauge
parameter, an off-block \(Y\) produces block-diagonal output or a
different pair. Since \(F^{-1}\) is pair-block diagonal,

% DISPLAY a.18
% SOURCE tr(F^(-1)L)=0.                                         (12)
\begin{gather*}
\begin{aligned}
& \operatorname{tr} (F^{-1}L)=0.
\end{aligned}\tag{12}
\end{gather*}

The remaining correction -K \(O_{C}\) is quadratic in \(Y\). In fact the
absolutely convergent logarithm series gives, for the stated \(\kappa\)
and \(\sigma\),

% DISPLAY a.19
% SOURCE |log(J/J_0)| <=8 q s^2.                                 (13)
\begin{gather*}
\begin{aligned}
& \vert \log (J/J_{0})\vert  \le 8 q s^{2}.
\end{aligned}\tag{13}
\end{gather*}

For verification:
\(\Vert F^{-1}L\Vert \le 2s/(1-\kappa ), \Vert F^{-1}K O_{C}\Vert \le 4s^{2}/(1-\kappa )\).
The linear trace vanishes, and the sum of powers \(\ge 2\) is bounded by
\(q e^{2}/[2(1-e)]\), where \(e=(2s+4s^{2})/(1-\kappa )\). Together with
the quadratic trace this is \(<8q s^{2}\).

\subsection{4. Intrinsic projectors and transition
law}\label{a-intrinsic-projectors-and-transition-law}

For an ordered orthogonal family of projectors \(P_{a}\) of ranks
\(n_{a}\) define the gauge-invariant function

% DISPLAY a.20
% SOURCE Phi_X(P)=sum_(a,i) n_a^(-1) (Tr(P_a X_i))^2.
\begin{gather*}
\begin{aligned}
& \Phi _{X}(P)=\sum _{a,i} n_{a}^{-1} (\operatorname{Tr} (P_{a} X_{i}))^{2}.
\end{aligned}\notag
\end{gather*}

Stationarity under off-block variations is precisely \(B_{Z}^{*}Y=0\).
At such a point its Hessian in a gauge parameter \(\xi\) is

% DISPLAY a.21
% SOURCE Hess Phi_X[xi,xi]
% SOURCE   =2||O_C xi||^2-2<B_Z xi,B_X xi>
% SOURCE   =-2<xi,R M xi>.                                      (14)
\begin{gather*}
\begin{aligned}
& \operatorname{Hess}  \Phi _{X}[\xi ,\xi ] \\
& =2\Vert O_{C} \xi \Vert ^{2}-2\langle B_{Z} \xi ,B_{X} \xi\rangle \\
& =-2\langle \xi ,R M \xi\rangle .
\end{aligned}\tag{14}
\end{gather*}

Thus RM is symmetric. In the normalized orbit variable \(\eta =R \xi\),
the negative half-Hessian is \(M R^{-1}\), which by (9) lies between
\(0.9696 I\) and \(1.0304\) I. Every such stationary decomposition is a
strict local maximum with a nonsingular transverse Hessian. The implicit
function theorem therefore gives a unique smooth nearby projector
branch, equivariant under SU(N), with no internal eigenvalue gap
assumption.

On overlap of charts describing the same branch, the ordered \(P_{a}\)
agree. Their block frames differ by \(K\) and, when the labels are
changed consistently, by a permutation of equal-rank blocks. Normal
frames differ by orthogonal maps on \(N_{Z}\). These are exact unitary
bundle changes on bosonic fibers and the fermionic module. \(J\) is
invariant; no extra nonunit half-density amplitude appears. All
expressions (3)-(14) and the kinetic identity below intertwine under
these changes. One may prove them in local frames without introducing
scalar angular IMS cutoffs.

This is a \textbf{same-branch} transition theorem. Comparing different
nested partitions, proving that their selectors agree with all
commuting-tree choices, and uniformly transporting every decision-tree
cutoff off the commuting cone remain separate obligations. Local
nonsingularity alone does not prove any of those global claims.

\subsubsection{4.1 Smooth fast projection on the same
branch}\label{a-smooth-fast-projection-on-the-same-branch}

The frozen pair matrices \(G_{ab}, K_{ab}\) and \(D_{ab}\) of the
general fast-energy lemma are obtained at \(Y=0\) from (15)-(16) and the
exact quadratic potential. Their bosonic frequencies and fermionic
spectral gaps are bounded below by fixed positive multiples of
\(r_{ab}\) under (1). Functional calculus therefore gives a smooth
positive Gaussian and a smooth negative-energy occupied subspace,
including at every internal eigenvalue collision. The resulting rank-one
fast-vacuum projection \(P\) is gauge equivariant. Local vacuum frames
may have a nonflat scalar Berry connection; none is declared parallel
here.

A same-branch chart change acts on every finite pair matrix by its
actual orthogonal/color/spin unitary representation. Spectral functional
calculus and the positive Gaussian construction commute with that
action. Hence \(P'=V P V^{*}\) exactly. The same is true after a smooth
invariant radial cutoff, for example with
\(h=(\sum _{a<b}r_{ab}^{-2})^{-1/2}\) and
\(\chi (\vert Y\vert /(\sigma  h))\), followed by fiberwise
normalization. There is no angular scalar partition in this
construction. Since \(r_{*}/\sqrt{k(k-1)/2}\le h\le r_{*}\), the cutoff
tail bounds differ only by fixed N-dependent constants.

For the uncut product \(U\), an internal derivative in \(A_{a}\)
differentiates only the pair factors incident to a. For \(U_{\sigma}\)
the derivative also differentiates its global scalar normalization; this
term is bounded by the same derivative norms, and its difference from
the uncut normalization is exponentially small in the indicated tail
scale. The resolvent formula for their spectral projections and the
Gaussian square-root formula give norm bounds \(C_{N,j} r_{ab}^{-j}\)
for each fixed number \(j\) of such derivatives, after the natural fast
dilation. These follow from external pair gaps, never from an internal
Hamiltonian resolvent. Mixed derivatives have the corresponding products
of inverse incident scales. Global phase choices are unnecessary because
the projection and its induced connection are intrinsic.

\subsection{5. Exact all-vector kinetic plus Gauss
form}\label{a-exact-all-vector-kinetic-plus-gauss-form}

Let \(p_{H}\) be the slice coordinate momentum stack, with the normal
connection in its center component and ordinary internal and fast
components. Let \(S_{f}\) denote the off-block fermionic Gauss
generator, in the same sign convention as \(O^{*}p=S_{f}\). Residual
K-equivariance remains imposed. The orthogonal tangential projection of
the orbit is

% DISPLAY a.22
% SOURCE O_T=S g_S^(-1) S^*O.
\begin{gather*}
\begin{aligned}
& O_{T}=S g_{S}^{-1} S^{*}O.
\end{aligned}\notag
\end{gather*}

For a raw equivariant section \(\phi\) define

% DISPLAY a.23
% SOURCE D phi=W^(1/2) M^(-*)[S_f phi-O^*S g_S^(-1)p_H phi].      (15)
\begin{gather*}
\begin{aligned}
& D \phi =W^{1/2} M^{-*}[S_{f} \phi -O^{*}S g_{S}^{-1}p_{H} \phi ].
\end{aligned}\tag{15}
\end{gather*}

Then the exact ambient reduced kinetic form, in coefficient-one units,
is

% DISPLAY a.24
% SOURCE t_raw[phi]=integral J {
% SOURCE    <p_H phi,g_S^(-1)p_H phi>+||D phi||^2 }.              (16)
\begin{gather*}
\begin{aligned}
& t_{\mathrm{raw}}[\phi ]=\int  J \{ \\
& \langle p_{H} \phi ,g_{S}^{-1}p_{H} \phi\rangle +\Vert D \phi \Vert ^{2} \}.
\end{aligned}\tag{16}
\end{gather*}

Proof: write an ambient covector as its slice-tangent projection plus
\(N \beta\). Its tangential restriction is \(p_{H}\); the Gauss equation
fixes \(M^{*} \beta =S_{f}-O^{*}S g_{S}^{-1}p_{H}\). Tangent and normal
parts are orthogonal, and
\(\Vert N \beta \Vert ^{2}=\langle \beta ,W \beta\rangle\). This proves
(16) for arbitrary covectors, arbitrary spinor values, and unrestricted
momenta. It is a first-order identity, not substitution into a squared
differential operator.

For the flat measure put \(d=\sqrt{J}\) and \(\psi =d \phi\). The exact
flattened formula is obtained from (15)-(16) by

% DISPLAY a.25
% SOURCE p_H -> p_H+i d_H log d.                                 (17)
\begin{gather*}
\begin{aligned}
& p_{H} \to  p_{H}+i d_{H} \log  d.
\end{aligned}\tag{17}
\end{gather*}

This fixes every ordering. Equivalently integration of the real density
cross terms gives the unshifted quadratic stack plus the scalar
multiplication

% DISPLAY a.26
% SOURCE V_d=d^(-1) div_H(a d_H d),                              (18)
\begin{gather*}
\begin{aligned}
& V_{d}=d^{-1} \operatorname{div} _{H}(a d_{H} d),
\end{aligned}\tag{18}
\end{gather*}

where a is the complete scalar principal co-metric, including the
orbital part of \(D\). Divergence is that of the flat bundle measure.
The finite Hermitian spin terms have zero real cross pairing with the
purely imaginary scalar density shift. No separate spin connection is
appended. Formula (17) remains the definitive expression if a convention
for momenta is changed.

On a fixed-shape regular region \(r_{*}\ge \delta  \rho _{C}\) with
\(\rho _{C}^{2}=\sum _{a} n_{a}\vert z_{a}\vert ^{2}\), and a closed
smaller tube, \(V_{d}\) and all coefficient derivatives needed at any
fixed order are bounded by constants times their homogeneous powers. In
particular
\(\vert V_{d}\vert \le C_{N,\delta ,\kappa ,\sigma} \rho _{C}^{-2}\).
These are bounds from explicit smooth finite matrices and compact shape
sets; they do not use a full Hamiltonian inequality.

Combining the exact internal kinetic coefficient with the exact internal
potential and Yukawa pieces retains

% DISPLAY a.27
% SOURCE sum_a [t_(A_a)+V_a+Yukawa_a]=sum_a h_(n_a)
\begin{gather*}
\begin{aligned}
& \sum _{a} [t_{A_{a}}+V_{a}+\mathrm{Yukawa}_{a}]=\sum _{a} h_{n_{a}}
\end{aligned}\notag
\end{gather*}

with coefficient one, also on colored internal sections. The
averaged-supercharge nonnegativity of these internal forms is the usual
algebraic statement, not a physical-sector Hardy inequality or an
internal spectral gap.

Equation (16) does \textbf{not} justify adding a separately extracted
adapted fast kinetic matrix and the whole same Gauss square. The frozen
adapted kinetic matrix is already part of that square.

\subsection{6. Incident-edge reserve for unrestricted low internal
derivatives}\label{a-incident-edge-reserve-for-unrestricted-low-internal-derivatives}

Let U(Z,A) be the normalized product of the pairwise kinetic-aware
Gaussian and occupied fermionic determinant from the frozen-pair lemma.
Interpret it as an isometric injection of its vacuum line; no flat Berry
connection is assumed. On each pair its covariance obeys

% DISPLAY a.28
% SOURCE E_U |Y_ab|^2 <= C n_a n_b ell^3/r_ab,
\begin{gather*}
\begin{aligned}
& E_{U} \vert Y_{ab}\vert ^{2} \le  C n_{a} n_{b} \ell ^{3}/r_{ab},
\end{aligned}\notag
\end{gather*}

with an absolute \(C\) on the stated \(\kappa\) tube. Use a smooth
invariant cutoff strictly inside (2), normalize fiberwise, and call the
resulting injection \(U_{\sigma}\). At fixed finite \(N\) and
sufficiently large \(r_{*}\), the same covariance bound holds with twice
\(C\); cutoff errors are bounded by polynomial factors times
\(\exp (-c \sigma ^{2} r_{*}^{3}/\ell ^{3})\), uniformly in unrelated A
and \(r_{\max}\). Derivatives can have fixed shape-dependent prefactors
if expressed in \(\rho _{C}\) units.

The internal orbital map in (15) is exactly

% DISPLAY a.29
% SOURCE O_(T,A)=O_A=P_I B_Y.                                    (19)
\begin{gather*}
\begin{aligned}
& O_{T,A}=O_{A}=P_{I} B_{Y}.
\end{aligned}\tag{19}
\end{gather*}

There is no Y-independent internal orbital derivative and no
contribution from an unrelated pair. For an internal covector
\(p=(p_{a})\), the ab component satisfies

% DISPLAY a.30
% SOURCE |(O_A^*p)_ab| <= C |Y_ab| (|p_a|+|p_b|).                 (20)
\begin{gather*}
\begin{aligned}
& \vert (O_{A}^{*}p)_{ab}\vert  \le  C \vert Y_{ab}\vert  (\vert p_{a}\vert +\vert p_{b}\vert ).
\end{aligned}\tag{20}
\end{gather*}

Write \(M=F+P\). The elementary bound

% DISPLAY a.31
% SOURCE ||P F^(-1)|| <=(2s+4s^2)/(1-kappa)<0.021
\begin{gather*}
\begin{aligned}
& \Vert P F^{-1}\Vert  \le (2s+4s^{2})/(1-\kappa )<0.021
\end{aligned}\notag
\end{gather*}

shows that \(M^{-*}=(I+F^{-*}P^{*})^{-1}F^{-*}\) has a uniformly bounded
left factor. \(W^{1/2}\) is uniformly bounded too. Thus (20) gives the
pointwise estimate

% DISPLAY a.32
% SOURCE ||W^(1/2)M^(-*)O_A^*p||^2
% SOURCE     <=C sum_(a<b) |Y_ab|^2/r_ab^2 (|p_a|^2+|p_b|^2).    (21)
\begin{gather*}
\begin{aligned}
& \Vert W^{1/2}M^{-*}O_{A}^{*}p\Vert ^{2} \\
& \le C \sum _{a<b} \vert Y_{ab}\vert ^{2}/r_{ab}^{2} (\vert p_{a}\vert ^{2}+\vert p_{b}\vert ^{2}).
\end{aligned}\tag{21}
\end{gather*}

This is valid for every internal covector, not just bounded internal
momenta. Compressing only its coefficients in \(U_{\sigma}\) yields

% DISPLAY a.33
% SOURCE ||D_A(U_sigma nabla_A f)||^2
% SOURCE     <=sum_a w_a |nabla_(A_a) f|^2,
% SOURCE w_a=C sum_(b!=a) n_a n_b ell^3 r_ab^(-3).                (22)
\begin{gather*}
\begin{aligned}
& \Vert D_{A}(U_{\sigma} \nabla _{A} f)\Vert ^{2} \\
& \le \sum _{a} w_{a} \vert \nabla _{A_{a}} f\vert ^{2}, \\
& w_{a}=C \sum _{b\ne a} n_{a} n_{b} \ell ^{3} r_{ab}^{-3}.
\end{aligned}\tag{22}
\end{gather*}

The derivatives here are the covariant derivatives on the possibly
nonflat vacuum line. The expression in (22) denotes the part where the
derivative acts on \(f\); derivatives acting on \(U_{\sigma}\) are
separate coefficient sources. Estimate (22) records the
incident-weighted low-leg coefficient bound. Its use in the mixed form
must follow the complete-charge and postcompletion estimates of Appendix
I, Sections \(4\) and 5.2-5.4; it does not authorize retaining or
spending a second copy of the original Gauss square.

Most importantly, the loss is incident-edge weighted. For an ordinary
untwisted internal derivative, both
\(t_{A_{a}}\le h_{a}+C_{n_{a}} \ell ^{-3} \alpha _{a}\) and
\(V_{a}\le h_{a}+C_{n_{a}} \ell ^{-3} \alpha _{a}\) follow from the
finite-Clifford Yukawa bound. Independently of that form comparison, (1)
gives

% DISPLAY a.34
% SOURCE w_a alpha_a <=C kappa sum_(b!=a)n_a n_b ell^3 r_ab^(-2). (23)
\begin{gather*}
\begin{aligned}
& w_{a} \alpha _{a} \le C \kappa  \sum _{b\ne a}n_{a} n_{b} \ell ^{3} r_{ab}^{-2}.
\end{aligned}\tag{23}
\end{gather*}

The weighted derivative error has a direct payment that does not require
a new weighted full-supercharge theorem. Use the internal-radial gauge
constructed in source Section \(4.2\) in Appendix \(F\): each pair
connection obeys \(\vert a_{ab}\vert \le C_{N} \kappa ^{2}/r_{ab}\), and
its \(A_{a}\) component vanishes unless that pair is incident to a. In
the resulting identification with the flat \(A=0\) center line,

% DISPLAY a.35
% SOURCE sum_a w_a ||nabla_(A_a)f||^2
% SOURCE    <=2 sum_a w_a t_(A_a)[f]
% SOURCE       +C_N kappa^4 ell^3 r_*^(-3)
% SOURCE                      integral sum_(a<b)r_ab^(-2)|f|^2
% SOURCE    <=2 sum_a w_a h_a[f]
% SOURCE       +C_N [kappa+kappa^4 ell^3 r_*^(-3)]
% SOURCE                      integral sum_(a<b)r_ab^(-2)|f|^2.   (26)
\begin{gather*}
\begin{aligned}
& \sum _{a} w_{a} \Vert \nabla _{A_{a}}f\Vert ^{2} \\
& \le 2 \sum _{a} w_{a} t_{A_{a}}[f] \\
& +C_{N} \kappa ^{4} \ell ^{3} r_{*}^{-3} \\
& \int  \sum _{a<b}r_{ab}^{-2}\vert f\vert ^{2} \\
& \le 2 \sum _{a} w_{a} h_{a}[f] \\
& +C_{N} [\kappa +\kappa ^{4} \ell ^{3} r_{*}^{-3}] \\
& \int  \sum _{a<b}r_{ab}^{-2}\vert f\vert ^{2}.
\end{aligned}\tag{26}
\end{gather*}

All rank-dependent Clifford and endpoint factors are included in
\(C_{N}\). The first inequality is the elementary
\(\vert \partial  f+i a f\vert ^{2}\le 2\vert \partial  f\vert ^{2}+2\vert a\vert ^{2}\vert f\vert ^{2}\);
the second uses (23). The weights depend only on centers, so multiplying
the internal form creates no omitted internal weight derivative. At
large \(r_{*}\) their \(h_{a}\) coefficients vanish. The inverse-square
term is arbitrarily small in \(\kappa\) and has a shape-independent
pair-center Hardy payment. This argument pays only the small derivative
error (22); it does not replace the full-supercharge comparison needed
for the principal Berry-coupled slow form. A cruder replacement
\(w_{a}\le C r_{*}^{-3}\), followed by \(r_{*}^{-3} \alpha _{a}\), would
destroy this conclusion for a far extended block and is not used.

The variable induced center metric is also paid as a genuine
derivative-form pairing against a retained covariant center kinetic
square. Its correction is exactly

% DISPLAY a.36
% SOURCE g_C^(-1)-I_C=-K^*(I+K K^*)^(-1)K.
\begin{gather*}
\begin{aligned}
& g_{C}^{-1}-I_{C}=-K^{*}(I+K K^{*})^{-1}K.
\end{aligned}\notag
\end{gather*}

Its Gaussian coefficient norm is \(O_{N}(\ell ^{3} r_{*}^{-3})\), so the
squared cost in that direct pairing is \(O_{N}(\ell ^{6} r_{*}^{-6})\)
times the low center derivative form. This step involves only center
derivatives; it cannot create an internal kinetic loss. Derivatives of
\(U_{\sigma}\) supply further coefficient sources, which still require
their own full source inventory.
